% titel: Blatt 9, A12: Welche Sulfide fallen aus?
% Teil II: der durchgerechnete Weg.

\subsection*{Gegeben und gesucht}
$pH=0.5$ enthält $c(Ni^{2+})=0.15$, $c(Co^{2+})=0.10$, $c(Cd^{2+})=0.5$
mol$L^{-1}$. Gesättigt mit H$_2$S (0.1 mol$L^{-1}$).
$pK_{a1}=7$, $pK_{a2}=14$. $L(NiS)=3\cdot10^{-21}$, $L(CoS)=5\cdot10^{-22}$,
$L(CdS)=1\cdot10^{-28}$. Fällt NiS, CoS, CdS aus?

\subsection*{Ansatz}
Erst $[S^{2-}]$, dann je Ionenprodukt $Q$, vergleichen mit $L$.

\subsection*{Rechnung}
\[
  [H^+] = 10^{-0.5} \approx 0.316\ mol\,L^{-1},\qquad
  K_{a1}=10^{-7},\ K_{a2}=10^{-14}.
\]
\[
  [S^{2-}] = \frac{K_{a1}K_{a2}\,[\mathrm{H_2S}]}{[H^+]^2}
    = \frac{10^{-7}\,10^{-14}\,0.1}{ (0.316)^2}
    = \frac{10^{-22}}{10^{-1}} = 10^{-21}\ mol\,L^{-1}.
\]
Ionenprodukte (jeweils $[M^{2+}]\cdot[S^{2-}]$):
\begin{align*}
  Q_{\text{Ni}} &= 0.15\cdot10^{-21} = 1.5\cdot10^{-22} < L(NiS)=3\cdot10^{-21}
    \Rightarrow \text{keine Fällung},\\
  Q_{\text{Co}} &= 0.10\cdot10^{-21} = 1\cdot10^{-22} < L(CoS)=5\cdot10^{-22}
    \Rightarrow \text{keine Fällung},\\
  Q_{\text{Cd}} &= 0.5\cdot10^{-21} = 5\cdot10^{-22} \gg L(CdS)=1\cdot10^{-28}
    \Rightarrow \textbf{Fällung}.
\end{align*}

\begin{ergebnis}
Nur \textbf{CdS fällt aus}; NiS und CoS bleiben gelöst (Ionenprodukt $Q$
unter dem jeweiligen $L$).
\end{ergebnis}

\begin{kontrolle}
Das extrem kleine $L(CdS)=10^{-28}$ macht CdS das mit Abstand
schwerlöslichste; selbst die nur mässige $[S^{2-}]$ reicht für Fällung.
Die Reihenfolge $L(Sulfid)$: CdS $\ll$ CoS $\lesssim$ NiS erklärt das
Verhalten. Physikalisch plausibel.
\end{kontrolle}

\begin{merksatz}
Sulfide fällen selektiv: $[S^{2-}]$ aus $K_{a1}K_{a2}[\mathrm{H_2S}]/[H^+]^2$,
dann $Q=[M^{2+}][S^{2-}]$ mit $L$ vergleichen. $Q>L$ $\Rightarrow$
Fällung.
\end{merksatz}
